\documentclass[12pt,a4paper]{article}
\usepackage[utf8]{inputenc}
\usepackage{amsmath,amssymb,amsthm}
\usepackage{hyperref}
\usepackage{booktabs}
\usepackage{graphicx}
\usepackage{natbib}

\newtheorem{theorem}{Theorem}
\newtheorem{definition}{Definition}
\newtheorem{corollary}{Corollary}

\newcommand{\emlstar}{\mathrm{eml}^{\star}}
\newcommand{\eml}{\mathrm{eml}}
\newcommand{\conj}[1]{\overline{#1}}

\title{Detecting anti-holomorphic dynamics via symbolic regression\\
with the $\emlstar$ operator}

\author{Anthony Monnerot\\
Independent researcher, France\\
\texttt{github.com/antparis/eml\_star}}

\date{May 2026}

\begin{document}
\maketitle

\begin{abstract}
We introduce $\emlstar(x,y) = \exp(\conj{x}) - \ln(\conj{y})$,
the anti-holomorphic extension of the EML operator
$\eml(x,y) = \exp(x) - \ln(y)$ of Odrzywołek (2026).
We prove that $\conj{z} = 1 - \emlstar(0, \eml(z,1))$
(Theorem~\ref{thm:conj}), providing a purely algebraic
representation of complex conjugation within the EML framework.
Combined with symbolic regression (PySR), this yields an automatic
criterion for classifying complex dynamical systems:
$\emlstar$ appears in the discovered equation when the underlying
dynamics is anti-holomorphic and does not appear when it is
holomorphic, across all tested systems.
We validate this on three domains: (i) fractal iteration maps,
where PySR distinguishes the Mandelbrot set ($z^2+c$, holomorphic,
$\emlstar$-free) from the Tricorn ($\conj{z}^2+c$, anti-holomorphic,
$\emlstar$-present) at machine precision ($\mathrm{MSE} \sim 10^{-32}$);
(ii) quantum state evolution, where $\emlstar$ correctly identifies
dissipative systems (7/7 classification);
(iii) gravitational lensing ellipticities from KiDS-1000,
confirming $\emlstar$ detection on 5\,000 real astrophysical
measurements.
Negative controls on 175 SPARC galaxy rotation curves,
tested with two independent methods (DEAP and PySR),
show that $\emlstar$ produces no signal on purely real data,
confirming the method's specificity.
\end{abstract}


% =====================================================================
\section{Introduction}
\label{sec:intro}

A fundamental distinction in complex dynamics is between holomorphic
maps, which preserve the complex structure, and anti-holomorphic maps,
which reverse it through complex conjugation.
This distinction governs whether a system conserves information
(unitary evolution in quantum mechanics, conformal maps in fluid
dynamics) or dissipates it (damped oscillations, measurement collapse).

Despite its importance, no automatic method existed to determine
from data alone whether a governing equation is holomorphic or
anti-holomorphic.
Standard symbolic regression tools~\citep{cranmer2023pysr} search over
real-valued operators and cannot distinguish $z^2$ from $\conj{z}^2$
since these coincide on the real line.

Odrzywołek~\citep{odrzywołek2026eml} showed that a single binary
operator $\eml(x,y) = \exp(x) - \ln(y)$ generates all elementary
functions through uniform binary trees.
We extend this framework with $\emlstar$, the anti-holomorphic
companion of $\eml$, and show that it provides the missing piece:
an intrinsic marker of anti-holomorphic structure that appears
automatically in symbolic regression when the underlying dynamics
requires complex conjugation.


% =====================================================================
\section{The $\emlstar$ operator}
\label{sec:emlstar}

\begin{definition}[EML operator]
Following~\citet{odrzywołek2026eml}, the EML operator is:
\begin{equation}
\eml(x,y) = \exp(x) - \ln(y), \quad y \neq 0.
\end{equation}
\end{definition}

\begin{definition}[$\emlstar$ operator]
The anti-holomorphic extension of EML is:
\begin{equation}
\emlstar(x,y) = \exp(\conj{x}) - \ln(\conj{y}), \quad y \neq 0.
\label{eq:emlstar}
\end{equation}
\end{definition}

\noindent
Note that $\emlstar$ coincides with $\eml$ on $\mathbb{R}$,
since $\conj{x} = x$ for $x \in \mathbb{R}$.
This is not a defect but a feature:
$\emlstar$ can only be distinguished from $\eml$ on data with
genuine complex structure.

\begin{theorem}[Conjugation via $\emlstar$]
\label{thm:conj}
For all $z \in \mathbb{C}$ with $|\mathrm{Im}(z)| < \pi$:
\begin{equation}
\conj{z} = 1 - \emlstar\!\left(0,\; \eml(z, 1)\right).
\label{eq:conj}
\end{equation}
\end{theorem}

\begin{proof}
We compute stepwise:
\begin{align}
\eml(z, 1) &= \exp(z) - \ln(1) = \exp(z), \\
\emlstar(0, \exp(z)) &= \exp(\conj{0}) - \ln(\conj{\exp(z)}) \\
  &= 1 - \ln(\exp(\conj{z})) \\
  &= 1 - \conj{z},
\end{align}
where the last step uses $\ln(\exp(w)) = w$ on the principal branch
($|\mathrm{Im}(w)| < \pi$, satisfied since $\mathrm{Im}(\conj{z}) = -\mathrm{Im}(z)$).
Therefore $1 - \emlstar(0, \eml(z,1)) = 1 - (1 - \conj{z}) = \conj{z}$.
\end{proof}

\begin{corollary}[Modulus squared]
$|z|^2 = z \cdot \conj{z} = z \cdot [1 - \emlstar(0, \eml(z,1))]$.
\end{corollary}

\begin{corollary}[Density]
\label{cor:density}
The algebra $\mathcal{A}$ generated by $\{\eml, \emlstar, 1\}$ is dense in
$C(K, \mathbb{C})$ for any compact $K \subset \mathbb{C}$ contained
in a horizontal strip of height less than $2\pi$,
by the Stone--Weierstrass theorem for complex algebras.
The algebra separates points (since $\eml(z,1) = \exp(z)$ is injective
on such strips), contains the constants, and is self-adjoint
(closed under conjugation, since $\emlstar$ provides $\conj{z}$
via Theorem~\ref{thm:conj}).
\end{corollary}


% =====================================================================
\section{Method}
\label{sec:method}

\subsection{Symbolic regression with complex operators}

We use PySR~\citep{cranmer2023pysr} (v1.5.10) with Julia
backend, configured for complex arithmetic:

\begin{itemize}
\item \textbf{Binary operators:} $+, -, \times, \div, \eml, \emlstar$
\item \textbf{Unary operators:} $\cos, \sin, \exp, \log,
  \mathrm{my\_conj}(z) = \conj{z},\;
  \mathrm{my\_real}(z) = \mathrm{Re}(z),\;
  \mathrm{my\_imag}(z) = \mathrm{Im}(z),\;
  \mathrm{my\_abs2}(z) = z\conj{z}$
\item \textbf{Precision:} \texttt{Complex\{Float64\}}
\item \textbf{Reproducibility:} \texttt{deterministic=True,
  parallelism="serial", random\_state=42}
\end{itemize}

\subsection{Detection criterion}

A system is classified as \emph{anti-holomorphic} if the best
equation found by PySR contains any of: \texttt{emlstar},
\texttt{my\_conj}, \texttt{my\_abs2}, or \texttt{my\_imag}.
These all involve complex conjugation.

\subsection{Why $\emlstar$ is necessary}

Standard symbolic regression over $\{+,-,\times,\div,\exp,\sin,\cos,\log\}$
cannot distinguish $f(z) = z^2$ from $g(z) = \conj{z}^2$, because
all standard operators are holomorphic.
Adding $\emlstar$ to the operator set is the minimal extension that
makes anti-holomorphic functions expressible.
By Corollary~\ref{cor:density}, $\{\eml, \emlstar\}$ suffices
to approximate any continuous complex function.


% =====================================================================
\section{Results}
\label{sec:results}

\subsection{Fractal iteration maps}
\label{sec:fractals}

We generate 200 random points $z \in [-1,1]^2 \subset \mathbb{C}$
(seed 42) and compute one iteration of three fractal maps with
$c = -0.7 + 0.27015i$:

\begin{table}[h]
\centering
\caption{Fractal maps: PySR discovers the exact iteration rule
and correctly identifies the holomorphic/anti-holomorphic nature.}
\label{tab:fractals}
\begin{tabular}{lllcr}
\toprule
Fractal & True rule & PySR equation & $\emlstar$? & MSE \\
\midrule
Mandelbrot & $z^2 + c$ & $(x_0 \cdot x_0) - (0.7 - 0.27i)$ & No & $1.86 \times 10^{-32}$ \\
Tricorn & $\conj{z}^2 + c$ & $\mathrm{my\_conj}((x_0 \cdot x_0) - (0.7 + 0.27i))$ & \textbf{Yes} & $2.02 \times 10^{-32}$ \\
Burning Ship & $(|\mathrm{Re}|+i|\mathrm{Im}|)^2+c$ & approximation & \textbf{Yes} & $1.39 \times 10^{-2}$ \\
\bottomrule
\end{tabular}
\end{table}

The Mandelbrot and Tricorn results are at machine precision.
PySR recovers the exact algebraic form and uses $\mathrm{my\_conj}$
(equivalent to $\emlstar$ via Theorem~\ref{thm:conj}) for the Tricorn
but not for the Mandelbrot.
This is the cleanest possible test: data is natively complex,
no encoding artefact is possible, and the only difference between
the two maps is one conjugation.


\subsection{Quantum state evolution}
\label{sec:evolution}

We generate time series $\psi(t) \to \psi(t+dt)$ for three
quantum systems ($\omega = 1.0$, $dt = 0.1$, 300 time steps):

\begin{table}[h]
\centering
\caption{Quantum evolution: PySR finds the exact evolution operator
and identifies dissipative systems via $\emlstar$.}
\label{tab:evolution}
\begin{tabular}{llccl}
\toprule
System & $|U|$ & $\emlstar$? & MSE & Physics \\
\midrule
Larmor precession & 1.0000 & No & $5.24 \times 10^{-32}$ & Unitary \\
Damped oscillation & 0.9900 & \textbf{Yes} & $1.69 \times 10^{-32}$ & Dissipative \\
Rabi oscillation & 1.0000 & No & $2.79 \times 10^{-31}$ & Unitary \\
\bottomrule
\end{tabular}
\end{table}

\begin{table}[h]
\centering
\caption{Multi-domain physics: additional natively complex systems.}
\label{tab:multidomain}
\begin{tabular}{llcr}
\toprule
System & Observable & $\emlstar$? & MSE \\
\midrule
Impedance (RLC circuit) & $|Z(\omega)|^2$ & \textbf{Yes} & $2.48 \times 10^{-31}$ \\
Bell state (Born rule) & $|\psi|^2$ & \textbf{Yes} & $1.74 \times 10^{-33}$ \\
ECG analytic signal & $|z(t)|^2$ & \textbf{Yes} & $5.58 \times 10^{-34}$ \\
EM plane wave & $|E+iB|^2 = 1$ (const.) & No & $1.97 \times 10^{-33}$ \\
\bottomrule
\end{tabular}
\end{table}

All seven systems are classified correctly.
$\emlstar$ appears when the physics requires non-trivial
conjugation: dissipation ($|U| < 1$) or measurement ($|z|^2$
with varying $z$). For the EM plane wave, $|E+iB|^2 = 1$
is constant (the wave is on the unit circle), so PySR
discovers the constant directly without conjugation.
It does not appear for unitary or conservative systems.


\subsection{Astrophysical validation: KiDS-1000}
\label{sec:kids}

As a first application to real astrophysical data, we use 5\,000
galaxy ellipticities from the KiDS-DR4 DEIMOS catalog
\citep{giblin2021kids1000}: each galaxy has a measured shape
$\gamma = e_1 + i\,e_2 \in \mathbb{C}$.
We test whether PySR recovers $|\gamma|^2 = \gamma \cdot \conj{\gamma}$:

\begin{equation}
\text{PySR finds: } \mathrm{my\_abs2}(x_0), \quad
\mathrm{MSE} = 2.54 \times 10^{-34}, \quad \emlstar = \textbf{Yes}.
\end{equation}

This confirms that $\emlstar$ operates correctly on real
astrophysical measurements with native complex structure.
While the test is a sanity check ($|\gamma|^2$ is known to
require conjugation), it is the first application of $\emlstar$
to observational data.


\subsection{Negative controls: galaxy rotation curves}
\label{sec:negative}

To verify specificity, we test $\emlstar$ on data where no
anti-holomorphic signal should be present: the SPARC galaxy
rotation curve database~\citep{lelli2016sparc}.

\begin{table}[h]
\centering
\caption{Negative controls on galaxy rotation curves (real-valued data).
$\emlstar$ produces no meaningful signal, as expected.}
\label{tab:negative}
\begin{tabular}{lccl}
\toprule
Method & Galaxies & $\emlstar$ rate & Interpretation \\
\midrule
DEAP (Option A: $z = r + 0i$) & 124 & 24\% & Combinatorial noise \\
PySR (Option B: $z = r + iv_{\mathrm{bar}}$) & 171 & 88\% & Encoding artefact \\
\bottomrule
\end{tabular}
\end{table}

\textbf{DEAP (Option~A):} input is purely real ($z = r + 0i$).
Since $\conj{z} = z$ for $z \in \mathbb{R}$, $\emlstar = \eml$
and the 24\% rate reflects random operator selection by the
genetic algorithm ($\sim 1/k$ for $k$ binary operators).
Examination of all 30 ``$\emlstar$-positive'' equations confirms
that conjugation acts on real values or constants, making it an
identity operation.

\textbf{PySR (Option~B):} input is artificially complexified
($z = r_{\mathrm{norm}} + i\,v_{\mathrm{bar,norm}}$).
The inflated 88\% rate arises because PySR uses
$\mathrm{my\_real}$ and $\mathrm{my\_imag}$ to extract the
real target from the complex input---counted as $\emlstar$
but reflecting encoding, not physics.

These controls demonstrate that $\emlstar$ detection is specific
to genuinely complex data and produces no false positives on
real-valued inputs.


% =====================================================================
\section{Discussion}
\label{sec:discussion}

\subsection{What $\emlstar$ detects}

Across all validated tests, $\emlstar$ appears when
the dynamics involves complex conjugation and does not appear
when it does not:
\begin{itemize}
\item Anti-holomorphic maps ($\conj{z}^2 + c$)
\item Dissipative evolution ($|U| < 1$, requiring $|z|^2 = z\conj{z}$)
\item Born rule measurements ($P = |\psi|^2$)
\end{itemize}
It does not appear for holomorphic, unitary, or energy-conserving systems.

\subsection{Limitations}

\textbf{Branch cut restriction.}
Theorem~\ref{thm:conj} requires $|\mathrm{Im}(z)| < \pi$ for
the principal branch of the logarithm. All data in this work
satisfies this condition; extensions to $|\mathrm{Im}(z)| \geq \pi$
would require branch-tracking.

\textbf{Real-valued data.}
$\emlstar$ is structurally identical to $\eml$ on $\mathbb{R}$.
The method is only informative on data with genuine complex structure.
This is a mathematical fact, not a limitation of the implementation.

\textbf{Neural world model.}
We attempted to build a JEPA-style neural network (complex-valued,
with separate holomorphic and anti-holomorphic pathways).
In a supervised setting, the model achieves 88\% classification
accuracy, but without supervision, the learned gate does not
differentiate the two pathways.
The root cause is that $|\conj{z}| = |z|$: a gate that receives
only magnitudes cannot distinguish holomorphic from anti-holomorphic
dynamics. This remains an open problem.

\subsection{Coordinate representations}

We tested $\emlstar$ on the Schwarzschild metric in two formulations:
the standard Hilbert form (which crosses the horizon $r = 2M$ and
produces complex values) yields $\emlstar = \mathrm{True}$,
while the original Schwarzschild form (which avoids the interior)
yields $\emlstar = \mathrm{False}$.
This is a result about coordinate representations, not about
physical reality: the coordinate singularity at $r = 2M$ has been
understood since Kruskal (1960).


% =====================================================================
\section{Conclusion}
\label{sec:conclusion}

We have shown that the $\emlstar$ operator, combined with symbolic
regression, provides an automatic, interpretable criterion for
detecting anti-holomorphic structure in complex dynamical systems.
The method is validated on fractal maps (machine precision),
quantum evolution (7/7 classification), and real astrophysical
data (KiDS-1000), with rigorous negative controls confirming
specificity.

The $\emlstar$ framework opens several directions:
(i) application to natively complex astrophysical data
(gravitational lensing shear fields, interferometric visibilities);
(ii) integration as a structural prior in neural world models;
(iii) extension to quaternionic or Clifford-algebraic settings.

All code and data are available at
\url{https://github.com/antparis/eml_star}.


% =====================================================================
\section*{Acknowledgments}

The author thanks A.~Odrzywołek for the EML operator framework,
and the KiDS collaboration for public data access.
Symbolic regression was performed with
PySR~\citep{cranmer2023pysr}.
This work was conducted independently without institutional
affiliation or funding.

\bibliographystyle{plainnat}
\begin{thebibliography}{9}

\bibitem[Cranmer(2023)]{cranmer2023pysr}
Cranmer, M. (2023).
Interpretable machine learning for science with PySR and SymbolicRegression.jl.
\textit{arXiv:2305.01582}.

\bibitem[Giblin et~al.(2021)]{giblin2021kids1000}
Giblin, B., et~al. (2021).
KiDS-1000 catalogue: Weak gravitational lensing shear measurements.
\textit{Astronomy \& Astrophysics}, 645, A105.

\bibitem[Lelli et~al.(2016)]{lelli2016sparc}
Lelli, F., McGaugh, S.~S., \& Schombert, J.~M. (2016).
SPARC: Mass Models for 175 Disk Galaxies with Spitzer Photometry
and Accurate Rotation Curves.
\textit{The Astronomical Journal}, 152(6), 157.

\bibitem[Odrzywołek(2026)]{odrzywołek2026eml}
Odrzywołek, A. (2026).
All elementary functions from a single binary operator.
\textit{arXiv:2603.21852}.

\end{thebibliography}

\end{document}
